\begin{trapbox}
\begin{description}[style=nextline, leftmargin=2.8em, labelsep=0.5em, itemsep=0.55em]
\item[\textbf{\textcolor{CTrap}{易错一：轴点写成 $0\times\infty$}}]
在 $P=(\pm a,0)$ 时切线竖直，$k_l$ 不存在；
在 $P=(0,\pm b)$ 时，$OP$ 竖直，$k_{OP}$ 不存在。
\textbf{纠正：} 乘积在这些点没有定义，不能规定 $0\times\infty=0$，
也不能把邻近点的极限直接当成原点处的乘积。
改用切线方程 $x=\pm a$ 或 $y=\pm b$。

\item[\textbf{\textcolor{CTrap}{易错二：少写负号，或写成 $1-e^2$}}]
对标准横轴椭圆 $0<e<1$，必有
$\boxed{-1<k_{OP}k_l=e^2-1<0}$。
\textbf{纠正：} 从切点式记住
$k_l=-\dfrac{b^2x_0}{a^2y_0}$；乘积必定为负。

\item[\textbf{\textcolor{CTrap}{易错三：把横轴椭圆公式硬套竖轴椭圆}}]
若方程改成
$\dfrac{x^2}{b^2}+\dfrac{y^2}{a^2}=1\ (a>b>0)$，
其斜率乘积实际为 $-a^2/b^2=-1/(1-e^2)$，
\textbf{不是} $e^2-1$。
\textbf{纠正：} 先辨认 $x^2$、$y^2$ 的分母，
横轴椭圆才使用本文主公式。

\item[\textbf{\textcolor{CTrap}{易错四：把中心 $O$ 偷换成焦点 $F$}}]
本结论说的是\textbf{椭圆中心与切点}之间的斜率 $k_{OP}$。
若是 $FP$，一般不能直接代入。
\textbf{纠正：} 先确认第一个点为 $O(0,0)$。

\item[\textbf{\textcolor{CTrap}{易错五：误认椭圆切线恒垂直于 $OP$}}]
圆的切线垂直半径，而非圆椭圆上一般不成立。
\textbf{纠正：} 圆可看成 $a=b$，这时乘积才是 $-1$；
对 $a>b$ 的椭圆，乘积在 $(-1,0)$ 内，不等于 $-1$。

\item[\textbf{\textcolor{CTrap}{易错六：将双曲线关系误记为 $e^2+1$}}]
对标准横轴双曲线
$\dfrac{x^2}{a^2}-\dfrac{y^2}{b^2}=1$，
同样设 $OP$ 与切线斜率存在，求导可得
$k_{OP}k_l=\dfrac{b^2}{a^2}=e^2-1>0$
（双曲线 $e^2=1+b^2/a^2$）。
\textbf{纠正：} 形式虽然也可写 $e^2-1$，
但\emph{椭圆为负、双曲线为正}；绝不能记作 $e^2+1$。

\item[\textbf{\textcolor{CTrap}{易错七：示意点未验证在曲线上}}]
例如对 $a=4,b=3$，$P=(2,1.5)$ 实际满足
$2^2/16+1.5^2/9=1/2\ne1$，不是椭圆上的切点。
\textbf{纠正：} 做题先代回方程检查；可以用
$P=(2,3\sqrt3/2)$ 作为真实例点。

\item[\textbf{\textcolor{CTrap}{易错八：已知切线斜率便唯一确定切点}}]
对关于中心对称的椭圆，$P$ 与 $-P$ 的切线平行，斜率相同。
\textbf{纠正：} 斜率条件通常只确定 $OP$ 的\emph{方向}，
求具体切点还要结合象限或其他条件。
\end{description}
\end{trapbox}
